\section{Problem-Solving Agents：以 MLAgentBench 为例}

\subsection{把“会聊”变成“会做”}
讲者展示了 problem-solving agent 的代表方向：在真实机器学习任务中完成端到端循环，包括读任务、改代码、跑实验、分析日志、迭代策略。

\begin{figure}[H]
\centering
\includegraphics[width=0.93\textwidth]{frame-06-groupchat.jpg}
\caption{复杂任务规划与求解：从对话到可执行 pipeline}
\end{figure}
\noindent\footnotesize 画面证据：字幕约在 00:10:02--00:13:10，讲者讨论 benchmark 任务求解、bash 执行和自动评测。\normalsize

\begin{importantbox}{Benchmark 价值：让“能力讨论”变成“可重复实验”}
没有 benchmark，agent 改进常停留在 demo 叙事。  
有 benchmark 后，才能比较不同模型、不同策略、不同工具链的真实增益，并发现失败分布。
\end{importantbox}

\subsection{评测目标与失败分解}
\begin{table}[H]
\centering
\begin{tabular}{p{2.8cm}p{4.0cm}p{4.8cm}}
\toprule
\textbf{评测维度} & \textbf{可观测信号} & \textbf{典型失败} \\
\midrule
任务完成度 & 指标是否达标、提交是否有效 & 目标漂移、任务误解 \\
执行可靠性 & 命令成功率、重试次数、超时率 & 环境依赖错误、工具调用失败 \\
策略质量 & 迭代效率、搜索覆盖、反思有效性 & 盲目试错、局部最优 \\
安全与约束 & 越权行为、异常副作用、审计可追踪性 & 未授权操作、不可回滚 \\
\bottomrule
\end{tabular}
\caption{Problem-solving agent 的评测不只看“最后分数”，还要看执行轨迹质量}
\end{table}

\begin{knowledgebox}{为什么 cyber tasks 会成为关键测试场景}
网络安全任务天然具备高反馈密度（成功/失败明确）与高操作风险（误操作代价高）。  
这类任务可以同时检验 agent 的执行能力、约束遵循能力和恢复能力。
\end{knowledgebox}

\begin{warningbox}{“刷分”并不等于真实能力提升}
若 benchmark 分布过窄，agent 可能学会特定套路而非通用策略。  
讲者强调应持续扩展任务分布与评测维度，避免出现“高分但不可迁移”的假进步。
\end{warningbox}

\subsection{本章小结}
Problem-solving agent 的关键进展来自“可执行 + 可评测 + 可复盘”。MLAgentBench 类工作真正推动的是研究范式，而不仅是单次榜单成绩。


